% page 205 du fac-simile (image p-204 du scan)
% bloc 167.6 x 263.6 mm ; marge gauche 34.4 mm
\pgc{12.700mm}{0.000mm}{
\l{}
\l{}
\l{}
\l{}
\l{\cel{55}197.}
\l{}
\l{\sou{gluo}. Materio viskoza, de qua strato extensita inter du objekti,}
\l{\cel{3}igas li inter-adherar. - EFI.}
\l{}
\l{\sou{guglar}. (\sou{netrans.})\cel{2}Bruisar quale la vino qua ekiras ek la kolo}
\l{\cel{3}di botelo. - DFIS.}
\l{}
\l{\sou{glumo}. (\sou{bot.}) Envelopilo florala di singla ek la spiketi qui}
\l{\cel{3}konstitucas la spiko di la gramini. (En la linguo ne-ciencala :}
\l{\cel{3}\sou{brano}). - EF.}
\l{}
\l{\sou{glutar}. (\sou{trans.}) Decensigar en la guturo. - II. (\sou{metaf}.) Desapar-}
\l{\cel{3}igar subite (en abismo, en abiso). - EFIRS.}
\l{}
\l{\sou{gluteo}. (\sou{anat.}) Parto mi-sferatra, karnoza, qua formacas singla}
\l{\cel{3}latero di la regiono dopa di la korpo di homo o di animalo.-EL.}
\l{}
\l{\sou{gluteino}.\cel{2}Materio viskoza qua restas en la farino de la cereali,}
\l{\cel{3}kande onu deprenis de ici la amilo.}
\l{}
\l{\sou{gluteno}. Substanco tenaca qua glutinas, interligas ad ensemblo,}
\l{\cel{3}la parti dividita di la korpi solida. - DEFIS.}
\l{}
\l{\sou{glutinar.(trans.})\cel{2}Adherigar per gluo.- EF.}
\l{}
\l{\sou{gneiso}. (\sou{mineral.})\cel{3}Roko ek feldspato, mikao, e c., ma di naturo}
\l{\cel{3}skistatra. - DEFIRS.}
\l{}
\l{\sou{gnomo}. (\sou{mitol.}) Mikra genio qua regnas tero e la kontenajo di ica.}
\l{\cel{3}- DEFIRS.}
\l{}
\l{\sou{gnomiko}. (\sou{antique}) Skriptisto Greka qua versigis sentenci etikala}
\l{\cel{3}o morala. - DEFIS.}
\l{}
\l{\sou{gnomono}. Stango vertikala di qua la ombro, projektate sur plano,}
\l{\cel{3}igas savar, per sua longeso, la alteso di Suno, ed indikas la}
\l{\cel{3}kloko per sua situeso. - DEFIRS.}
\l{}
\l{\sou{gnoso}. Doktrino (religio-hereziala) segun qua la enti spiritala,}
\l{\cel{3}ekirinte de la sino di Deo, retroeniros lu. - DEFIS.}
\l{}
\l{\sou{gnostiko}. Adepto di gnoso. - DEFIRS.}
\l{}
\l{\sou{gnuo}. (\sou{zool}.) Speco di granda antilopo di Afrika. L.\cel{1}\sou{catoblepas}}
\l{\cel{3}\sou{gnu}.- DEFIRS.}
\l{}
\l{\sou{gobio}. (\sou{zool.}) Mikra fisho manjebla, de la genero \textquotedbl{}ciprini\textquotedbl{}, qua}
\l{\cel{3}vivas trupe apud la fundo\cel{2}di la riveri. - L.\cel{1}\sou{cyprinus gobio}.}
\l{\cel{3}- DSL.}
\l{}
\l{\sou{gobleto}. Drinko-vazo, alta, cilindratra, sen anso, ed ordinare}
\l{\cel{3}sen pedo. - EFS.}
\l{}
\l{\sou{godrono}. I. (\sou{tekn.}) Muluro ye seciono ovala, borde di la repasto-}
\l{\cel{3}vazaro arjenta. - II. (\sou{stofo}).Plisuro ye seciono cirkla en la}
\l{\cel{3}kolumi plisita, en lejaboti. en la mansheti. - EF.}
}
